\chapter{Estado del arte}
\label{cap:estado-arte}

\section{Introducción y delimitación del problema}

La detección de música generada mediante inteligencia artificial se ha ido consolidando como un problema propio dentro del análisis de audio. En trabajos recientes, esta tarea aparece bajo la denominación \textit{Synthetic Song Detection} (SSD), entendida como la clasificación entre canciones producidas por humanos y canciones generadas de extremo a extremo mediante sistemas de IA \cite{rahman2024sonics}.

Esta delimitación es importante porque existen tareas cercanas que conviene distinguir. La \textit{Singing Voice Deepfake Detection} se centra en detectar voces cantadas sintéticas o convertidas, a menudo dentro de mezclas donde el acompañamiento puede proceder de música real \cite{zang2023singfake,zang2024ctrsvdd}. La detección de canciones sintéticas de extremo a extremo, en cambio, considera que pueden ser artificiales la voz, la instrumentación, la letra, el arreglo y la producción global. SSD también se diferencia de la detección de voz hablada sintética, la detección de audio sintético general, la identificación de letras generadas mediante IA y la atribución del modelo generador.

El interés de SSD va más allá de obtener una decisión automática sobre un archivo de audio. La música es una señal con estructura temporal, patrones tímbricos, organización rítmica y relaciones de larga duración. Por ello, los detectores pueden apoyarse tanto en indicios locales de la señal como en propiedades musicales más amplias. Esta combinación explica que el campo incluya desde enfoques clásicos basados en características acústicas hasta modelos profundos que intentan aprovechar representaciones aprendidas y relaciones entre fragmentos.

\section{Conjuntos de datos para detección de música sintética}

En el contexto de SSD, SONICS aporta un conjunto de datos específico para detectar canciones generadas de extremo a extremo. Sus autores señalan que muchos corpus previos se centraban en voz cantada sintética, mientras que los sistemas actuales pueden generar también la instrumentación, el arreglo y la producción global \cite{rahman2024sonics}. Esta diferencia motiva la creación de conjuntos específicos para canciones completas y de protocolos capaces de considerar relaciones temporales más amplias.

FakeMusicCaps aborda un problema próximo desde otra perspectiva. El conjunto de datos se construye a partir de MusicCaps, regenerando pares música-texto mediante varios modelos \textit{text-to-music}. Además de la detección binaria, el trabajo estudia la atribución del modelo generador \cite{comanducci2024fakemusiccaps}. FakeMusicCaps constituye un ejemplo de corpus específico para música generada mediante IA y muestra cómo la detección puede verse condicionada por los generadores incluidos en el conjunto de datos.

AIME ocupa una posición distinta. La tarjeta oficial del conjunto de datos lo presenta como \textit{AI Music Evaluation Dataset}, con música generada por distintos modelos y ejemplos procedentes de MTG-Jamendo, junto con metadatos como el modelo de origen y una descripción basada en etiquetas \cite{discoethAIME}. Su finalidad original está vinculada a la evaluación de modelos de generación musical y preferencias humanas, por lo que su utilización en detección requiere definir y documentar un protocolo experimental propio \cite{groetschla2025benchmarking}.

La reutilización de AIME para detección cuenta al menos con un antecedente verificado: \textit{Fusion Segment Transformer}, identificado como preprint de 2026, evalúa su propuesta en SONICS y AIME \cite{kim2026fusionsegmenttransformer}. Esta evidencia es todavía limitada, aunque muestra que el corpus puede emplearse en trabajos de detección cuando se explicita el protocolo experimental.

\section{Enfoques clásicos con características acústicas}

Los enfoques clásicos de clasificación de audio suelen partir de características acústicas calculadas manualmente. Estas representaciones resumen propiedades de la señal antes de aplicar un clasificador, lo que permite construir sistemas relativamente simples, reproducibles y de bajo coste computacional. En este contexto, los MFCC representan información espectral de corto plazo mediante vectores compactos.

Los MFCC se han utilizado en tareas relacionadas de detección de audio sintético, especialmente en sistemas de anti-\textit{spoofing} de voz. El precedente aplicado citado en este capítulo es Novoselov et al., que emplean MFCC dentro del reto ASVspoof 2015 y comparan clasificadores clásicos con enfoques más complejos \cite{novoselov2015stcantispoofing}. Aunque este dominio no coincide con SSD, sirve como antecedente metodológico para entender su uso como representación acústica compacta en problemas de detección.

La SVM es un clasificador clásico para tareas supervisadas, especialmente adecuado cuando se trabaja con vectores de características fijas. Su formulación de margen máximo y el uso de kernels permiten construir fronteras de decisión no lineales sin requerir arquitecturas profundas \cite{cortes1995support}. La combinación MFCC + SVM constituye un baseline clásico reproducible y controlable, basado en características acústicas compactas.

La principal limitación de este tipo de aproximación es que resume la señal en descriptores compactos. Esto puede ser suficiente para capturar ciertos patrones espectrales, pero deja menos margen para modelar relaciones temporales largas, estructura musical o dependencias entre secciones de una canción. Estas limitaciones han favorecido el estudio de representaciones profundas capaces de modelar información acústica y temporal más compleja.

\section{Enfoques profundos y representaciones preentrenadas}

El desarrollo reciente de SSD se ha desplazado hacia modelos profundos capaces de aprender representaciones más ricas a partir del audio. SONICS propone SpecTTTra, una arquitectura orientada a capturar dependencias temporales largas con menor coste que algunos modelos convolucionales o transformadores de referencia \cite{rahman2024sonics}. Esta línea pone de manifiesto que la detección de canciones completas depende tanto de indicios locales como de la forma en que la música se organiza a lo largo del tiempo.

Otra línea dentro de los enfoques profundos consiste en reutilizar conocimiento aprendido previamente, en lugar de entrenar toda la arquitectura desde cero para el problema concreto. Los encoders de audio pueden entrenarse sobre corpus amplios y después emplearse como extractores de vectores para tareas de clasificación.

Dentro de esta línea se sitúa MERT, entrenado mediante aprendizaje autosupervisado a gran escala sobre audio musical \cite{li2024mert}. Su propuesta combina objetivos acústicos y de alto nivel durante el preentrenamiento para capturar información reutilizable. Los vectores resultantes se evaluaron en problemas de recuperación y comprensión de información musical, lo que sitúa a MERT como un encoder especializado en este dominio.

\section{Segmentación y evaluación a nivel de canción}

La segmentación aparece de forma natural en detección musical porque las canciones pueden tener duraciones superiores a las que muchos modelos procesan directamente. Trabajar con fragmentos permite adaptar el audio a modelos con entradas de duración fija, acotar el coste de cada inferencia y obtener varias observaciones de una misma pieza. La decisión a nivel de canción requiere, además, una estrategia de agregación adecuada.

Segment Transformer divide el audio en segmentos, extrae información de cada uno mediante modelos preentrenados o autosupervisados y aprende relaciones entre fragmentos para mejorar la detección a escala de canción \cite{kim2025segmenttransformer}. La propuesta muestra la utilidad de separar el análisis local de los fragmentos y el modelado estructural posterior, especialmente cuando el audio completo es demasiado largo para tratarse de una sola vez.

El preprint \textit{Fusion Segment Transformer} amplía esta idea mediante una fusión guiada de información de contenido y estructura, y proporciona una evidencia directa de evaluación sobre AIME dentro de un trabajo de detección \cite{kim2026fusionsegmenttransformer}. Esta línea refuerza la relevancia de los enfoques por segmentos en problemas de detección musical.

En tareas de clasificación musical, la forma de dividir el audio y agregar predicciones puede afectar a las medidas de rendimiento. Nasrullah y Zhao muestran, en clasificación de artistas, que la longitud de los fragmentos, el esquema de particionado y la agregación a nivel de canción influyen en la evaluación \cite{nasrullah2019musicartistclassification}. Aunque la tarea es distinta, el resultado es relevante para diseñar protocolos de segmentación que eviten la mezcla de información relacionada entre particiones.

Los trabajos específicos de SSD también refuerzan esta idea. SONICS destaca la importancia de las dependencias de larga duración, mientras que Segment Transformer modela explícitamente relaciones entre segmentos \cite{rahman2024sonics,kim2025segmenttransformer}. Esto sugiere que un detector puede operar sobre fragmentos y, al mismo tiempo, necesitar una estrategia de agregación si se desea emitir una decisión sobre una canción completa.

\section{Generalización, robustez y limitaciones actuales}

Una dificultad recurrente en SSD es la generalización fuera del escenario de entrenamiento. Afchar et al. muestran que obtener una medida alta en una prueba cerrada resulta insuficiente para considerar resuelto el problema, ya que la robustez frente a manipulaciones del audio y la generalización a generadores no vistos siguen siendo cuestiones críticas \cite{afchar2025challenges}. Esta advertencia es especialmente importante cuando los detectores pueden aprender atajos acústicos asociados a un conjunto de datos, a un códec o a un proceso de producción concreto.

Los cambios de dominio también aparecen en trabajos adyacentes de SVDD. SingFake identifica dificultades con cantantes no vistos, códecs de comunicación, idiomas y contextos musicales \cite{zang2023singfake}. CtrSVDD, por su parte, destaca que la selección de características y la separación entre métodos de falsificación vistos y no vistos afectan al comportamiento del detector \cite{zang2024ctrsvdd}. Aunque pertenecen a una tarea distinta, estos resultados muestran la fragilidad de los detectores cuando cambian las condiciones de evaluación.

Los preprints recientes de 2026 insisten en esa misma dirección desde la detección de música generada mediante IA. MusicDET plantea explícitamente la detección ante generadores no vistos mediante un enfoque \textit{zero-shot} \cite{han2026musicdet}. Broadcast Monitoring analiza un escenario donde la música aparece en fragmentos cortos o enmascarada por habla, y muestra que las condiciones de emisión pueden degradar modelos entrenados para contextos más limpios \cite{lopezayala2026broadcast}. ArtifactNet, también como preprint, centra la atención en la robustez frente a códecs y en detectores ligeros basados en residuos forenses \cite{oh2026artifactnet}.

En conjunto, estas aportaciones muestran que las medidas obtenidas en un corpus cerrado deben interpretarse con cautela. La duración de los fragmentos, el idioma, el artista, el género, la frecuencia de muestreo, la compresión y los artefactos de producción pueden influir en la decisión del modelo. Por tanto, una comparación experimental rigurosa debe distinguir entre rendimiento dentro del conjunto de datos evaluado y capacidad de generalización a condiciones externas.

\section{Conclusiones}

Synthetic Song Detection es una tarea reciente y todavía en desarrollo. Aunque ya existen conjuntos de datos y arquitecturas específicas para música generada de extremo a extremo, todavía falta un enfoque consolidado que pueda considerarse solución general. Los resultados en escenarios cerrados son útiles para comparar modelos, pero ofrecen garantías limitadas ante nuevos generadores, transformaciones de audio o cambios de dominio.

La revisión realizada permite justificar la comparación entre un enfoque basado en MFCC y SVM y un enfoque profundo basado en MERT congelado y SVM bajo condiciones comunes de evaluación. El primero aporta una referencia reproducible basada en características manuales, mientras que el segundo permite evaluar representaciones profundas preentrenadas orientadas al dominio musical.

Por último, AIME debe presentarse con precisión: es un corpus documentado procedente de la evaluación de modelos de generación musical, reutilizado aquí para detección binaria tras una auditoría propia. La clase IA incluye doce generadores, lo que evita depender de un único generador dentro del corpus, aunque no garantiza generalización a generadores o datasets no vistos.
